\section{Data}
\label{sec:data}

Four image populations enter this work, all built on Euclid Q1 VIS
images~\citep{euclidq1overview, euclidq1data}: simulated lenses with matched negatives, forming
the labelled source domain; expert-graded real lenses, whose labels are used only at
evaluation; an unlabelled pool of real Euclid galaxies; and objects graded X (not a lens) by expert
inspection, used as hard negatives. Figure~\ref{fig:setup} gives the structure of the splits; this section
gives the reasons.

\subsection{Arc-on-real simulations}

The source domain is simulated, and only the arc in it is simulated. Each positive is a
real Euclid VIS cutout of a galaxy --- the deflector --- with a ray-traced arc
painted on; each matched negative is the \emph{same} galaxy with no arc. The deflectors are
galaxies from the Q1 catalogue whose stored values satisfy $\log M_* \geq \DeflMassFloor{}$
and $\DeflZmin{} \leq z_{\rm phot} \leq \DeflZmax{}$~\citep{euclidq1phz}. We generated
30\,000 such pairs from 19\,167 of them, roughly 4\,182 of which can
produce an arc with Einstein radius $\theta_E \ge 0.75''$. Painting onto real cutouts
retains the real line-of-sight objects, noise and point-spread-function variation,
confining the sim-to-real gap to the arc; this is standard
practice~\citep{holismokesXI, lines2025discoveryengineC, unions2025}, although the lensed
sources here are parametric models rather than the real HST source galaxies that
Euclid Collaboration: Lines et al.~\citep{lines2025discoveryengineC} and \citet{unions2025} lens
(Section~\ref{sec:sim-pipeline}). The matched bare deflector is
the point of the construction: a positive and its negative differ in the arc and nothing
else, so the host galaxy carries no class information (Figure~\ref{fig:cutouts}).

An earlier fully synthetic generator lacked that property. Its simulated lenses had no
deflector light while the real negatives had bright centres, and a classifier reached
AUROC 0.995 on per-image brightness alone. Replacing it with the arc-on-real construction
--- real deflector light, raw flux, softened S\'ersic parameters, supersampled rendering
--- brought the residual morphology gap between the classes to approximately 0.84 in
AUROC. The shortcut had been invisible in every metric we were watching at the time.

\subsection{The simulation pipeline}
\label{sec:sim-pipeline}

The deflector's velocity dispersion $\sigma_v$ was assigned by abundance matching to the
function below; the Einstein radius $\theta_E$ followed from $\sigma_v$ under a singular
isothermal ellipsoid (SIE); and the source redshift $z_s$ was drawn so that $\theta_E$
fell in an acceptable range given the deflector's photometric redshift
$z_l$~\citep{euclidq1phz}. The background source is a
\emph{parametric model}, not a real HST image: sources were
drawn with GalSim~\citep{rowe2015galsim} as parametric models (single S\'ersic or
bulge-plus-disk) fitted to COSMOS F814W galaxies from its $I < \simSourceSample$ sample, the
training sample and fits described by~\citet{mandelbaum2014great3}, through
\texttt{COSMOSCatalog} with \texttt{gal\_type="\simSourceGalType"},
passing no catalogue exclusion arguments, and kept if their half-light radius lay between
$\simSourceHlrMin''$ and $\simSourceHlrMax''$. The source was ray-traced through the SIE with
\texttt{lenstronomy}~\citep{birrer2018lenstronomy, birrer2021lenstronomy}, convolved with a Gaussian PSF of FWHM $0.16''$,
scaled so that its brightest arc pixel reached a signal-to-noise ratio drawn uniformly
between \simArcSnrMin{} and \simArcSnrMax{} against the cutout's background, and added to
the cutout; no source magnitude cut was applied. Poisson noise was added to the
arc alone, the cutout already carrying real Euclid noise. The output is 60 shards of
30\,000 lenses and 30\,000 matched non-lenses at $64 \times 64$ pixels and $0.1''$/pixel;
realised Einstein radii have 10th, 50th and 90th percentiles $0.77''$, $0.92''$, $1.4''$.

Velocity dispersions were abundance-matched to the SDSS DR5 early-type (E/S0) velocity
dispersion function of~\citet{choi2007vdf}, a modified Schechter function deconvolved for
measurement error. The implementation uses
\begin{equation}
  \frac{\mathrm{d}n}{\mathrm{d}\sigma} \;\propto\;
  \left(\frac{\sigma}{\sigma_*}\right)^{\alpha}
  \exp\!\left[-\left(\frac{\sigma}{\sigma_*}\right)^{\beta}\right] \frac{1}{\sigma}
  \label{eq:vdf}
\end{equation}
with $\sigma_* = 161$\,km\,s$^{-1}$, $\alpha = 2.32$ and $\beta = 2.67$, on a uniform grid
of 6\,000 points over $\sigma \in [80, 400]$\,km\,s$^{-1}$, normalised numerically so the
amplitude does not enter. The dispersions were sorted and assigned to deflectors by
stellar-mass rank --- a rank-to-rank map rather than a fit --- under a fixed seed. Two
assumptions are ours rather than part of that measurement: the function describes central
stellar velocity dispersions, which we use directly as the SIE velocity dispersion; and it
is measured at $z \lesssim 0.1$, which we apply to every deflector without redshift
evolution.

Every image enters the encoder through one fixed transformation applied before
augmentation: a per-image clip at the 0th and 99.7th percentiles, a rescale to $[0,1]$,
and an $\operatorname{arcsinh}$ stretch at $q = 1$, bounded by
$\operatorname{arcsinh}(1) \approx 0.881$, resized to $224 \times 224$. This is the
\texttt{fits\_to\_jpg\_like\_scaling} recipe as implemented in the \texttt{galaxy-datasets}
package, reproduced from that source code rather than from a paper, applied once rather than twice;
re-stretching an augmented image would skew the per-image percentiles. At training time a
random positional shift was applied to every class equally. Random centroid shifts are
established practice and not a detail: training and testing on perfectly centred images
drops a baseline ResNet lens finder's AUROC from $0.9913$ to $0.8332$ and its
true-positive rate at zero false positives to $0\%$, the network otherwise learning spatial
offset as a determinant feature~\citep{holismokesXI}. Their shift is $\pm 5$ pixels on
HSC, about $\pm 0.84''$; the $\pm 3$ pixels used here is this work's own choice, which at
Euclid VIS's $0.1''$/pixel is $0.3''$ --- a smaller shift in angular terms, roughly a
third of the median Einstein radius, and inside the encoder's own training-augmentation
range. Applying it equally is what stops it becoming a class cue.

\subsection{Real lenses and the unlabelled pool}

The positive evaluation class is 489 real strong lenses from the Euclid Q1 Discovery
Engine catalogue, 246 of catalogue grade A and 243 of grade B~\citep{discoveryengineA},
taken with their cutouts from the catalogue's data record~\citep{discoveryenginecat}.
Experts graded each image on a response scale of A (confident lens), B (probable lens), C
(possible lens, with features that can be explained without lensing) and X (not a lens);
the scores of about ten experts per image are averaged with a correction for each expert's
optimism, and the catalogue grades are cuts on that aggregate, grade A above 2 and grade B
above 1.5~\citep{discoveryengineA}. The positive class is catalogue grades A and B,
following the benchmark~\citep{lines2025simstosky}; catalogue grade C falls below the B cut
and is excluded. The same survey team's ensemble study treats grade C candidates as
non-lenses for calibration, on the grounds that they are typically not lensed
systems~\citep{discoveryengineE}. The catalogue lists 497 grade A/B systems; eight are absent from the processed set for reasons the download and cutout logs record individually and none of which is a selection rule --- three appear in none of the released archives, three have no FITS file in the archive member that holds them, and two failed on read --- leaving 489. All 489 were used regardless of which archive stored them, the class being
defined by expert grade rather than by any downstream modelling outcome.

The unlabelled pool is approximately 225\,000 real Euclid VIS cutouts, cut from the
Q1 MER mosaics, of objects selected from the MER catalogue by VIS PSF-fit flux between
\FluxMin{} and \FluxMax{}~$\mu$Jy ($I_E \approx \MagFaint{}$ to $\MagBright{}$ in PSF
magnitude, which for resolved galaxies reads fainter than the total
magnitude)~\citep{euclidq1mer}. The same query required a VIS detection and positive Y, J
and H template-fit fluxes, a point-like probability below 0.3, an extended probability above
0.5 and a spurious probability below 0.3, within four boxes covering the Q1 fields, and was
capped at 250\,000 rows with no ordering specified. Since the cap was reached and the query
specifies no ordering, this parent pool is the archive's first 250\,000 matches rather than a
random draw from all matching objects. Objects are assigned to roles by a hash of their object
ID, which depends on no object property, and the training, validation and target negatives are
then drawn with a fixed seed. Separately,
4\,455 objects graded X --- not a lens --- by expert inspection in version 0.0.1 of the
Discovery Engine catalogue~\citep{discoveryenginecat001}, drawn from the highest-scoring
candidates of its networks and visual inspection~\citep{discoveryengineA}, were included
among the source negatives as hard negatives; later versions publish grades A--C
only~\citep{discoveryenginecat}, and \GradeXnowC{} of them are graded C in version 0.0.3
and were kept as negatives. These are the arc mimics that produce real false positives, and
curating such a set is established practice~\citep{holismokesXI, astrovink}. They appear in training
only, and in no evaluation pool.

\subsection{Splits and fold rotation}

Assignments are fixed by a single splits file built deterministically under seed 42, the
sole source of truth for every membership. There are four roles.

\textbf{Source-train} is 24\,000 simulated positives against 24\,000 negatives: all 4\,455
grade-X hard negatives, 2\,881 unique bare deflectors and 16\,664 random unlabelled
galaxies. The simulated lenses were split 80/20 into train and validation \emph{grouped by
deflector identifier}, one deflector having been reused for up to ten lenses, so ungrouped
splitting would place near-identical images on both sides~\citep{astrovink}.
\textbf{Source-validation} is the held-out 6\,000 positives against 6\,000 negatives from
a disjoint partition, with no grade-X hard negatives. With no real lens labels available
for training, model selection cannot follow the supervised workflow and must rest on a
proxy~\citep{rlsbench}; the proxy available here is held-out source data. It is a weak one: \citet{musgrave2021realitycheck} find that
source accuracy screens out poor models but does not identify good ones --- low source
accuracy goes with low target accuracy, while high source accuracy does not imply high
target accuracy --- and Section~\ref{sec:eval} returns to it.

\textbf{Target} rotates: for fold $k$ it is the roughly 391 graded real lenses of the
\emph{other four} folds together with roughly 391 random real galaxies from a reserved
partition. The real lenses in the target pool are the graded lenses of the other folds,
and they enter without their labels --- they supply real-lens density to the adaptation
objective, and no label of theirs reaches any loss. The pool is balanced at approximately
1:1 by construction, which is a theoretical requirement rather than a convenience: when
source and target label distributions differ, marginal-alignment methods provably cannot
minimise target error~\citep{zhao2020gls, classimbalanceda}, and our source is balanced 1:1.
\citet{parul2024} obtained a working marginal-alignment result in an astronomical setting
on this construction. Deployment prevalence is not represented here; it is handled at
inference by threshold calibration, and because the primary metric depends only on score
ranking (Section~\ref{sec:eval}) balancing the target does not distort it. Target
negatives are drawn \emph{fresh per fold} and are disjoint across folds: they are part of
training, so sharing them would give the fold models overlapping training data and
correlate them, and the spread across folds would then understate the uncertainty of the
fold mean. In cross-validation, non-overlapping training partitions have been reported to
give less optimistic variance estimates than overlapping ones~\citep{irredundant_kfold}.

\textbf{Test} is per fold: that fold's roughly 98 held-out graded lenses --- the one place
a real lens label is used --- against a single \emph{shared} pool of \TestNegDOne{} random real
galaxies (\TestNegDrawn{} as drawn, less the \DOneRemoved{} excluded below), identical for every fold. Sharing is the deliberate opposite of the target
choice, for the opposite reason: the metric is a true-positive rate at a threshold set by
a false-positive rate on the negatives, so one shared pool holds that threshold fixed and
makes fold-to-fold variation reflect positive recovery rather than negative-tail noise.
The size also matches the established Q1 convention of order 30\,000 test
negatives~\citep{lines2025simstosky, astrovink}; there, ``random'' describes how the pool
was drawn, and each object is then individually inspected (graded non-lens internally, or
rejected by citizen science), not left unvetted. Ours are drawn from the parent pool by ID hash and seeded sampling only, passed
through the screening cut above and no further individual inspection, which is
the deployment regime a search actually faces. The benchmark draws its non-lens sample from a Q1
search population limited at $I_E < 22.5$~\citep{lines2025simstosky}; ours extends to $\MagFaint{}$ in PSF magnitude, and
\FracFaintTest{}\% of our test negatives are fainter than 22.5 in that quantity, so the two
pools are not matched in depth. Test negatives exclude any object sharing an
exact Euclid ID with a catalogued lens or a simulation host galaxy (\DOneRemoved{}
objects); values on the pre-registered pool are given in Appendix~\ref{app:prereg_pool}.

The 489 lenses were divided into five folds stratified by grade and grouped by object
identifier, so each fold holds a comparable A/B mixture and any duplicated physical lens
stays on one side. Four folds hold 98 test positives and one holds 97; each lens is tested
exactly once, and the partition into a test role and a target role is exact on every fold.

\subsection{Leakage audit}

The splits were checked in two independent layers. The builder's own audit enforces nine
groups of rules --- covering disjointness, class balance, partition provenance, fold
rotation, grade representation, the confinement of hard negatives to training, and the
exclusion of simulation deflectors from every negative role --- and refuses to write the
file unless all of them hold. They are listed in Appendix~\ref{sec:splitrules}. All pass
on the real data. A second layer is not a formality: an independent Q1 lens-finding
pipeline that also splits at catalogue level, for the same reason, still needed a
perceptual-hash follow-up to catch four near-duplicate pairs its catalogue-level split had
let through~\citep{astrovink}.

A separate verifier then re-derived the checks the builder cannot make, in the spirit of a
medical-imaging foundation-model survey whose treatment of leakage in adaptation pipelines
motivates this audit~\citep{domain_leakage}, and found two things. Perceptual
hashing of the 489 lenses flagged eight similar pairs, whose sky positions are all more
than 3\,900 arcseconds apart, so they are distinct systems sharing the canonical
bright-elliptical-plus-ring appearance and were recorded as inspected and cleared. A
cross-match of simulation deflectors against the raw negative pool found 20 simulation
hosts in it; the builder was changed to purge simulation-deflector identifiers before
drawing, and the verifier then reported none in any negative role. Both checks compared
identifiers as the files store them, and the deflector files and simulation shards store
them rounded. Matched by exact Euclid ID, rebuilt from the deflector files' coordinates,
\DeflAllNeg{} simulation hosts are among the negatives: the \DOneDeflNeg{} in the test pool
are excluded by the test-set definition above, and \DeflOtherNeg{} remain among the
training, validation and target negatives. The simulation shards carry no sky coordinates,
so a host present under a different identifier would not be caught by either check.

\begin{figure}[t]
\centering
\includegraphics[width=\textwidth]{F1_setup.pdf}
\caption{\textbf{The four data roles, the fold rotation, and which arm sees what.} Every
count is read from the splits file; the shared test negatives are counted after the
exclusion described in the text. \textbf{No real lens label enters a training loss}:
the graded lenses appear in the target pool as unlabelled images and in the test pool as
labels used at evaluation only. Target negatives are fresh per fold so the fold models
stay uncorrelated; test negatives are one shared pool so the FPR${}=10^{-3}$ threshold is
identical across folds. The $\sim$1:1 target balance follows from the label-shift argument
above, and it is also what caps the pool at 782 images: the available positives are the
other folds' graded lenses, and matching the source ratio fixes the negatives to the same
count. That 782-image pool is the exposure the memorisation diagnostic reported later
measures.}
\label{fig:setup}
\end{figure}

\begin{figure}[t]
\centering
\includegraphics[width=\textwidth]{F2_cutouts.pdf}
\caption{\textbf{One example from each image population.} All four panels use the same
stretch --- the per-image arcsinh the encoder receives --- so they are comparable as
inputs to the model rather than in absolute flux. The examples follow a declared rule
rather than being chosen by eye: the simulated lens sits at its shard's \emph{median}
Einstein radius, and the second panel is \textbf{the same real galaxy with no arc painted
on}; the real lens is the grade-A lens at the median expert score. At the median Einstein
radius the arc is faint, which is what a median rather than a hand-picked example looks
like. In a UMAP embedding of their own feature space, Euclid Collaboration: Lines et al.~\citep{lines2025simstosky} find that
real Q1 lenses tend to have fainter arcs and less complete Einstein rings than their
simulations, and attribute part of the sim-to-real gap to that; whether our simulations are
deficient in that same region we did not measure, and we note their finding without
asserting the connection. Image credit: ESA Euclid/Euclid Consortium/NASA/Q1-2025.}
\label{fig:cutouts}
\end{figure}
